\chapter{جمع‌بندی، نتیجه‌گیری و رهنمون‌هایی برای آینده}
هدف کلی محاسبات کوانتومی توزیع شده، این است تا بتواند مسائل محاسباتی پیچیده را با تقسیم و تجزیه آن‌ها به چند زیر سیستم کوانتومی حل نماید و بار محاسباتی را بین این زیر سیستم‌ها تقسیم نماید. در فصل اول معماری کلی یک سیستم کوانتومی توزیع شده نشان داده شد که از چندین زیر لایه تشکیل شده بود. در این رساله هدف ارائه یک سیستم کوانتومی توزیع شده در لایه‌ی دوم با رویکرد های زیر بود:
\begin{itemize}
\item
سیستم پیشنهادی بایستی بتواند، بار محاسباتیی را بین زیر سیستم‌های کوانتومی توزیع نماید.
\item
سیستم پیشنهادی در دو سطح تعداد دوربری‌های کوانتومی را کاهش می‌دهد، به طوری که در سطح اول کیوبیت‌های سیستم را بین چند زیر سیستم با هدف کاهش تعداد ارتباطات مورد نیاز توزیع نموده و در سطح دوم تعداد دوربری‌های کوانتومی مورد نیاز را کاهش می‌دهد.
\end{itemize}

